% Setting up document class and essential packages
\documentclass[12pt]{article}
\usepackage[utf8]{inputenc}
\usepackage[T1]{fontenc}
\usepackage{amsmath}
\usepackage{amsfonts}
\usepackage{geometry}
\geometry{a4paper, margin=1in}
\usepackage{parskip}
\usepackage{listings}
\usepackage{xcolor}

% Configuring listings for Python code
\lstset{
    language=Python,
    basicstyle=\ttfamily\small,
    keywordstyle=\color{blue},
    stringstyle=\color{red},
    commentstyle=\color{green!60!black},
    showstringspaces=false,
    breaklines=true,
    frame=single
}

% Configuring fonts (loaded last)
\usepackage{noto}

\begin{document}

\section*{Numerical Roe Flux Workflow for General Hyperbolic Conservation Laws}

This document outlines a numerical procedure to compute the Roe flux for a general 1D hyperbolic conservation law \(\mathbf{u}_t + \mathbf{f}(\mathbf{u})_x = 0\), where \(\mathbf{u} \in \mathbb{R}^m\) is the conserved variable vector and \(\mathbf{f}(\mathbf{u})\) is the flux. The approach is system-agnostic, using numerical approximations to handle arbitrary systems without requiring analytical forms of the equations.

\subsection*{Workflow}
\begin{enumerate}
    \item \textbf{Numerical Flux Form}: Compute the Roe flux:
    \[
    \mathbf{F}_{i+1/2} = \frac{1}{2} \left( \mathbf{f}(\mathbf{u}_L) + \mathbf{f}(\mathbf{u}_R) \right) - \frac{1}{2} \sum_{k=1}^m |\tilde{\lambda}_k| \tilde{\alpha}_k \tilde{\mathbf{r}}_k
    \]
    where \(\mathbf{u}_L, \mathbf{u}_R \in \mathbb{R}^m\) are left and right states, \(\tilde{\lambda}_k\), \(\tilde{\alpha}_k\), \(\tilde{\mathbf{r}}_k\) are Roe-averaged eigenvalues, wave strengths, and eigenvectors.

    \item \textbf{Numerical Jacobian}: Approximate the flux Jacobian \(\mathbf{A} = \frac{\partial \mathbf{f}}{\partial \mathbf{u}}\) at \(\mathbf{u} \approx (\mathbf{u}_L + \mathbf{u}_R)/2\) using finite differences:
    \[
    \mathbf{A}_{ij} \approx \frac{f_i(\mathbf{u} + \epsilon \mathbf{e}_j) - f_i(\mathbf{u} - \epsilon \mathbf{e}_j)}{2\epsilon}
    \]
    where \(\mathbf{e}_j\) is the \(j\)-th unit vector, \(\epsilon \approx 10^{-6}\).

    \item \textbf{Eigenstructure}: Compute eigenvalues \(\tilde{\lambda}_k\) and eigenvectors \(\tilde{\mathbf{r}}_k\) of \(\mathbf{A}\) using a numerical eigensolver (e.g., QR algorithm or NumPy).

    \item \textbf{Wave Strengths}: Solve for \(\tilde{\alpha}_k\) in \(\Delta \mathbf{u} = \mathbf{u}_R - \mathbf{u}_L = \sum_{k=1}^m \tilde{\alpha}_k \tilde{\mathbf{r}}_k\) by projecting:
    \[
    \tilde{\alpha}_k = (\tilde{\mathbf{l}}_k)^T \Delta \mathbf{u}
    \]
    where \(\tilde{\mathbf{l}}_k\) are left eigenvectors.

    \item \textbf{Assemble Flux}: Compute:
    \[
    \mathbf{F}_{i+1/2} = \frac{1}{2} \left( \mathbf{f}(\mathbf{u}_L) + \mathbf{f}(\mathbf{u}_R) \right) - \frac{1}{2} \sum_{k=1}^m |\tilde{\lambda}_k| \tilde{\alpha}_k \tilde{\mathbf{r}}_k
    \]

    \item \textbf{Implementation Notes}:
    \begin{itemize}
        \item Use numerical libraries (e.g., NumPy, SciPy) for matrix operations.
        \item Ensure \(\mathbf{f}(\mathbf{u})\) is provided as a function.
        \item Apply entropy fix: replace \(|\tilde{\lambda}_k|\) with \(\max(|\tilde{\lambda}_k|, \delta)\), \(\delta \approx 10^{-6}\).
        \item Test with small \(\epsilon\) for numerical stability.
    \end{itemize}
\end{enumerate}

\subsection*{Python Implementation}
The following Python code implements the numerical Roe flux for a general system, demonstrated with the p-System (\(\mathbf{u} = [v, w]^T\), \(\mathbf{f} = [w, -\sigma(v)]^T\), \(\sigma(v) = v^2\)).

\begin{lstlisting}
import numpy as np

def compute_roe_flux(u_L, u_R, flux_func, epsilon=1e-6, delta=1e-6):
    """
    Compute Roe flux for a general 1D hyperbolic conservation law.
    
    Parameters:
    u_L, u_R : ndarray, left and right state vectors
    flux_func : function, computes flux f(u)
    epsilon : float, perturbation for finite difference Jacobian
    delta : float, entropy fix parameter
    
    Returns:
    F : ndarray, Roe flux at interface
    """
    m = len(u_L)  # Number of components
    u_avg = 0.5 * (u_L + u_R)  # Simple average for Jacobian evaluation
    
    # Numerical Jacobian
    A = np.zeros((m, m))
    for j in range(m):
        u_plus = u_avg.copy()
        u_minus = u_avg.copy()
        u_plus[j] += epsilon
        u_minus[j] -= epsilon
        A[:, j] = (flux_func(u_plus) - flux_func(u_minus)) / (2 * epsilon)
    
    # Eigenstructure
    eigenvalues, R = np.linalg.eig(A)
    L = np.linalg.inv(R)  # Left eigenvectors
    
    # Wave strengths
    delta_u = u_R - u_L
    alpha = L @ delta_u
    
    # Entropy fix
    abs_eigenvalues = np.array([max(abs(lam), delta) for lam in eigenvalues])
    
    # Roe flux
    F_avg = 0.5 * (flux_func(u_L) + flux_func(u_R))
    correction = 0.5 * sum(abs_eigenvalues[k] * alpha[k] * R[:, k] for k in range(m))
    F = F_avg - correction
    
    return F

# p-System flux function
def p_system_flux(u, sigma_func=lambda v: v**2):
    """
    Flux function for p-System: u = [v, w], f = [w, p(v)], p(v) = -sigma(v)
    """
    v, w = u
    return np.array([w, -sigma_func(v)])

# Example usage for p-System
def main():
    # Example states
    u_L = np.array([1.0, 0.5])  # [v_L, w_L]
    u_R = np.array([0.8, 0.3])  # [v_R, w_R]
    
    # Compute Roe flux
    F = compute_roe_flux(u_L, u_R, p_system_flux)
    
    print("Left state u_L:", u_L)
    print("Right state u_R:", u_R)
    print("Roe flux F:", F)

if __name__ == "__main__":
    main()
\end{lstlisting}

This implementation is flexible, relying on numerical approximations and a provided flux function, making it applicable to any hyperbolic conservation law.

\end{document}